\chapter{Conclusion}
\label{ch:conclusion}

\section{Summary}
\label{sec:concl-summary}

This thesis builds the covariance input that the rank-inference procedure of
\citet{mogstad2024} requires when forecast losses are serially and cross-sectionally
dependent, and leaves the procedure itself unchanged. Newey--West standard errors on the
pairwise difference series, with a Bartlett kernel and the \citet{andrews1991}
data-dependent bandwidth, absorb serial dependence. Two constructions of the simulation
covariance handle cross-sectional dependence, the common shock that cancels in
differences. The first reconstructs the level covariance by classical multidimensional
scaling; it is calibrated by construction and repaired to positive semi-definiteness by
a nearest-PSD projection. The second, the direct estimator, estimates the difference
covariance, with a pairwise-of-pairs generalization for unbalanced panels. The two
coincide on a balanced panel, exactly so under a common bandwidth. On an unbalanced panel
only the direct estimator targets the covariance of the tested contrasts, because the
contrasts are then no longer differences of a single vector of means. A calibration
analysis explains the finite-sample difference between the two: the direct construction
studentizes by its own diagonal and stays calibrated, while the reconstruction
studentizes by the original standard errors and can lose calibration once its distance
matrix is far from Euclidean.

The Monte Carlo study confirms nominal coverage for both constructions and their
equivalence on clean panels, and isolates the crossover at which the direct construction
becomes the sharper of the two: non-stationarity combined with imbalance. Real data show
the same pattern. On the unbalanced U.S.\ SPF panels the reconstruction returns the
trivial set for every forecaster, while the direct construction separates three of the
five best-covered NGDP forecasters; on the near-balanced euro-area, tournament and
competition panels the two agree to within one rank. Across the three applications the
procedure behaves faithfully over the range of separability. It reports wide, mostly
trivial rank sets for macro forecasters that can hardly be distinguished, and a
partially separated top tier for tournament forecasters. For the methods submitted to
the M3 competition, ranked on the Paris temperature record, it separates twelve
competent methods from three systematic failures and pins the failures to the last three
ranks, but it names no winner among the competent twelve. In each case the width of the
confidence set measures how far the data can rank.

\section{Limitations}
\label{sec:concl-limitations}

Several caveats bound the results. Participation in the tournament is self-selected: the
panel consists of the twenty most active volunteers, who chose to take part in these
questions, so the ranking describes highly engaged forecasters on the questions they
selected, not a population ranking of skill. The competition panel is balanced, but its
effective sample is small. Its 306 observations come from seventeen forecast origins,
and the eighteen horizons forecast from one origin share the same misjudged level, so the
panel carries far less information than 306 independent comparisons would. The
tournament and competition applications also exercise mainly the cross-sectional side of
the construction: tournament questions have only a weak time order (by closing date),
and the competition's serial dependence is confined to the horizons within a block. In
all three applications, moreover, the loss differences the HAC layer sees are only mildly
persistent; the macro panels' serial correlation sits mainly in the individual errors.
The serial-dependence layer is therefore exercised most fully in simulation. The
competition application ranks \emph{methods}, as their authors implemented them for M3,
on a temperature record rather than an economic target. It is included because it shows
the procedure separating systematic failure from competence, not for its substantive
content. Its loss, sMAPE, is also not invariant to the origin of the measurement scale,
so the ranking is specific to the competition's scaling of the record. Finally, the whole
construction rests on the additive common-shock model; a genuinely multiplicative
interaction between the common shock and idiosyncratic skill would not cancel cleanly in
differences.

\section{Future work}
\label{sec:concl-future}

The most valuable extension is an application that is both informative and rich in both
kinds of dependence. Financial-analyst earnings forecasts fit this profile: a large,
unbalanced panel of forecasters repeatedly predicts an evolving target subject to common
market shocks. Publicly accessible analogues (crowd-sourced estimate platforms, ongoing
forecasting tournaments) are candidates. On the methodological side, the direct
construction could be extended to the $\tau$-best simulation variant. The calibration
diagnostic, namely the projection distance and the dispersion of the studentized draws,
could be developed into a formal test of when to distrust the reconstruction in favour of
the direct estimator.
